\chapter{צוות וממשל תאגידי}\label{ch:team}

חברה שמגדלת ירקות על גגות מנהלת בפועל שלושה עסקים במקביל: חקלאות, תפעול עירוני ותוכנה. לכן הצוות המייסד נבנה משלושה תחומי מומחיות משלימים, ולכל אחד מהם בעלים אחד שאחראי על תוצאה מדידה. פרק זה מציג את המייסדים, את תוכנית הגיוס, את המבנה הארגוני ואת מנגנוני הממשל שיעמדו לרשות המשקיעים.

\section{הצוות המייסד}

שלושת המייסדים (טבלה~\ref{tab:founders}) עובדים יחד על הפרויקט ומחזיקים בחלק העיקרי של ההון המניות. הרקע של כל אחד מהם מכסה את אחד משלושת סיכוני הליבה של העסק: ביצועי הגידול, אמינות התוכנה ויכולת הפריסה בשטח.

\begin{table}[htbp]
\centering\small
\rrstripes
\begin{tabularx}{\linewidth}{>{\raggedright}p{2.6cm}>{\raggedright}p{3.4cm}>{\raggedright\arraybackslash}X}
\rrtablehead \rrth{שם} & \rrth{תפקיד ורקע} & \rrth{אחריות מרכזית} \\
נועה ברק & מנכ"לית (\en{CEO}); אגרונומית, לשעבר מנהלת חממה & אסטרטגיה, גיוס הון, איכות הגידול ויחסי לקוחות מסעדות \\
עידו שחר & מנהל טכנולוגיות (\en{CTO}); מהנדס למידת מכונה ורובוטיקה & פיתוח \en{RootOS}, מודל ראיית העלה, תאום דיגיטלי ותשתית חיישנים \\
מאיה לוי & סמנכ"לית תפעול (\en{COO}); תפעול, לשעבר בחברת משלוחי מזון & פריסת גגות, שרשרת אספקה, לוגיסטיקה והסכמים עם ועדי בתים \\
\end{tabularx}
\caption{הצוות המייסד: תפקידים ואחריות.}
\label{tab:founders}
\end{table}

\medskip
שלושת התחומים נפגשים בנקודה אחת: גג שמניב \RRroofKg{} ק"ג בשנה במחיר ממוצע של \RRprice~\shekel{} לק"ג. נועה אחראית לכמות, עידו לאיכות הנתונים שמאפשרת לשמור עליה, ומאיה לכך שהתוצרת מגיעה בזמן ללקוח. התפקידים אינם חופפים, ולכן החלטה על גג חדש עוברת תמיד דרך שלושתם.

\begin{insightbox}[למה הצוות מתאים]
הסיכון הגדול ביותר בחקלאות עירונית הוא פער בין מעבדה לשטח. אגרונומית עם ניסיון בחממות מסחריות, מהנדס שבנה מערכות בקרה ומנהלת תפעול שהפעילה צי שליחים מכסים את שלושת הצדדים של הפער הזה כבר ביום הראשון.
\end{insightbox}

\section{תוכנית גיוס}

הצוות גדל בהתאם לקצב פריסת הגגות ולא לפניו. המטרה היא להגיע לקרוב ל-\RRheadcount{} עובדים בסוף 2028, כפי שמוצג בתרשים~\ref{fig:orgchart}, ולהמשיך לגדול בעיקר בתפקידי השטח, כי הם הנגזרת הישירה של מספר הגגות. מטה ההנהלה גדל לאט יותר, ולכן העלות ליחידה יורדת עם הצמיחה.

\begin{enumerate}
  \item \textbf{2027, שנת ההקמה.} שלושת המייסדים ומגדל עירוני ראשון. המתקינים והמהנדס מועסקים בקבלנות כדי לשמור על גמישות.
  \item \textbf{2028.} עד חמישה מגדלים עירוניים בסך הכול, מתקין ומהנדס שעוברים להעסקה ישירה, מפתחי \en{RootOS} ומנהל לקוחות \en{B2B} שכפוף למנכ"לית. הגיוסים נעשים לפני הגל השני של התקנות, כדי שהצוות יהיה מוכן כשהמערכות יעלו.
  \item \textbf{2029, אחרי סבב \en{Series A}.} ראש מכירות ושותפויות להפצת ה-\en{SaaS} לחוות חיצוניות, ותפקיד כספים מלא במקום מיקור חוץ חלקי.
  \item \textbf{2030 ואילך.} צוותי שטח אזוריים שמתרחבים לפי מספר הגגות, ומנהל איכות ובטיחות עם הכשרה בתחום עבודה בגובה.
\end{enumerate}

עלות המטה וההנהלה מופיעה בשורת ההוצאות הקבועות בטבלה~\ref{tab:pnl} ובפרק~\ref{ch:financials}. כל שכר ותפקיד שמתווספים נבחנים מול נקודת האיזון שם, ולכן הגיוס מתנהל לפי אבני דרך ולא לפי תקציב כללי.

\begin{figure}[htbp]
\centering
\resizebox{\linewidth}{!}{% Organisation at end of 2028 — mirrors data/headcount.csv (\RRheadcount people).
\begin{tikzpicture}[level distance=17mm,
  every node/.style={rounded corners=2mm,align=center,font=\small,
                     minimum width=2.6cm,minimum height=10mm},
  ceo/.style={fill=rrForest,text=white,font=\small\bfseries},
  vp/.style={fill=rrLeaf!45,draw=rrForest,line width=0.6pt},
  tm/.style={fill=rrCream,draw=rrStone,font=\scriptsize,minimum width=2.2cm},
  level 1/.style={sibling distance=37mm},
  level 2/.style={sibling distance=26mm},
  edge from parent/.style={draw=rrForest,line width=0.8pt}]
  \node[ceo] {מנכ"ל\\ \en{CEO}}
    child {node[vp] {סמנכ"ל תפעול\\ \en{COO}}
      child {node[tm] {5 מגדלים\\ עירוניים}}
      child {node[tm] {מתקין\\ + מהנדס}}}
    child {node[vp] {מנהל טכנולוגיות\\ \en{CTO}}
      child {node[tm] {2 מפתחי\\ \en{RootOS}}}}
    child {node[vp] {מנהל לקוחות\\ \en{B2B}}}
    child {node[vp] {כספים\\ (חלקי)}};
\end{tikzpicture}
}
\caption{המבנה הארגוני בסוף 2028 (כ-\RRheadcount{} עובדים). מערך התפעול הוא הגדול ביותר, כי הוא גדל יחד עם מספר הגגות.}
\label{fig:orgchart}
\end{figure}

\section{חבר יועצים}

כדי לכסות תחומים שאינם בליבת הצוות, יוקם חבר יועצים קטן, ללא שמות בשלב זה (טבלה~\ref{tab:advisors}). לכל יועץ תפקיד מוגדר, והוא נפגש עם ההנהלה אחת לרבעון.

\begin{table}[htbp]
\centering\small
\rrstripes
\begin{tabularx}{\linewidth}{>{\raggedright}p{3.6cm}>{\raggedright\arraybackslash}X}
\rrtablehead \rrth{תחום היועץ} & \rrth{התרומה לחברה} \\
מהנדס קונסטרוקציה & ביקורת על סקרי העומס לפני כל חוזה, מענה ישיר לסיכון המבני \\
אגרונום-חוקר בהידרופוניקה & אימות פרוטוקולי הגידול ושיטות המדידה של התפוקה \\
יועץ משפטי לנדל"ן ומקרקעין & נוסח החוזה מול ועדי הבתים, היתרים ואחריות \\
יועץ למסעדנות ושרשרת אספקה & חיבור ללקוחות \en{B2B} וקביעת מדיניות מחיר \\
יועץ \en{ESG} ומימון בר-קיימא & מסגרת הדיווח, גיוס ממשקיעי השפעה ועמידה בסטנדרטים \\
\end{tabularx}
\caption{חבר היועצים המתוכנן: תחומים ותרומה.}
\label{tab:advisors}
\end{table}

\section{ממשל תאגידי ודיווח \texorpdfstring{\en{ESG}}{ESG}}

הממשל בנוי לשקיפות כלפי משקיעים כבר מהסיבוב הראשון. הכללים פשוטים וניתנים לבדיקה.

\begin{itemize}
  \item \textbf{דירקטוריון.} שלושה חברים עד סבב \en{Series A}: נציג המייסדים, נציג המשקיעים ודירקטור בלתי תלוי. אחרי הסבב מתווסף חבר נוסף מטעם משקיעי הסבב.
  \item \textbf{ישיבות ודוחות.} ישיבת דירקטוריון אחת לרבעון, ודוח חודשי למשקיעים עם מספר הגגות, הכנסות, מזומן וסטייה מהתחזית.
  \item \textbf{החלטות שדורשות אישור.} חריגה מהתקציב השנתי, חוזה חריג מול ועד בית והחלטה על תוספת הון.
  \item \textbf{ניהול סיכונים.} רישום סיכונים מתעדכן אחת לרבעון, כמפורט בפרק~\ref{ch:risks}.
\end{itemize}

בדיווח \en{ESG} החברה תמדוד שלושה סוגי מדדים: סביבה (צריכת מים ליחידת תוצרת, שימוש בחשמל וקילומטרים של שינוע), חברה (גגות פעילים, מקומות עבודה והכנסה שמועברת לוועדי הבתים) וממשל תאגידי (מספר ישיבות, ניגודי עניינים ותיעוד החלטות). המדדים הללו אינם נבדקים במסגרת תוכנית זו ונקבעים כיעדים לשנה הראשונה.

\begin{insightbox}[המשמעות למשקיע]
צוות משלים, גיוס שקשור למספר הגגות וממשל שקוף מקטינים את הסיכון שתלוי באנשים, סיכון \en{R7} בטבלה~\ref{tab:risks}. המשקיע מקבל נתונים חודשיים ולא רק דוח שנתי.
\end{insightbox}
